\section{Additional results for permanent and polynomial inequalities}\label{app:additional-results}

This appendix gives exact counterexamples and supporting reductions related to Gyires's inequality. Each counterexample states the assertion that fails. None is a counterexample to Gyires's inequality for nonnegative doubly stochastic matrices. The results below have not been formalized in Lean.

\subsection{Definitions}\label{app:additional-1}

Let \(c_n=n!/n^n\), and let \(J_n\) have every entry equal to \(1/n\). A doubly stochastic matrix has nonnegative entries, and every row and column sum is one. For a real square matrix \(A\), put

\[
p_A(x)=\prod_{i=1}^n\sum_{j=1}^n a_{ij}x_j
       =\sum_{|\alpha|=n}p_\alpha x^\alpha,
\qquad N(p)=\sum_\alpha\alpha!p_\alpha^2,
\qquad \alpha!=\prod_j\alpha_j!.
\]

The permanent coefficient identity gives

\[
[x_1\cdots x_n]p_A=\per(A),\qquad
\langle p_A,p_B\rangle_F=\sum_\alpha\alpha!p_\alpha q_\alpha
=\per(AB^{\T}).
\]

Here \(q_\alpha\) are the coefficients of \(p_B\). The squared Fischer norm is \(N(p)=\langle p,p\rangle_F\). Thus the strong Gram inequality is

\begin{equation}
\per(A)^2\ge c_n N(p_A)
=c_n\per(AA^{\T}).
\label{eq:additional-1}
\end{equation}

For the original symmetrized denominator, write

\[
Q(A)=\frac{\per(AA^{\T})+\per(A^{\T}A)
+2\per(A^2)}4.
\]

Gyires's inequality is \(\per(A)^2\ge c_nQ(A)\). Applying \eqref{eq:additional-1} to both \(A\) and \(A^{\T}\) implies this inequality by permanent Cauchy--Schwarz.

A degree-\(m\) polynomial \(q\) in \(m\) variables is balanced when \(\partial_jq(\mathbf1)=q(\mathbf1)\) for every \(j\). For a polynomial with nonnegative coefficients, its capacity is

\[
\operatorname{cap}(q)=\inf_{x_j>0}\frac{q(x)}{x_1\cdots x_m}.
\]

At a positive diagonal capacity minimizer \(D\), the polynomial \(q(Dx)/\det D\) is balanced. Write \(H(q)=N(q(Dx))/(\det D)^2\) for its squared Fischer norm. When a finite minimizer does not exist, use the balanced polynomial limit under determinant-one capacity-minimizing scales. The examples that use \(H\) below give explicit balancing formulas. The notation \([x_j]p\) extracts exactly the first power of \(x_j\), leaving a polynomial in the other variables.

A polynomial is real stable if it has no zero when every variable has positive imaginary part. A nonnegative homogeneous stable polynomial is Lorentzian. In particular, each quadratic derivative of a degree-\(n\) Lorentzian polynomial has a Hessian with at most one positive eigenvalue. The examples below use this necessary condition. When that condition is used to establish Lorentzianity, the support condition is also stated. For a set \(S\subseteq\mathbb N^m\) of fixed coordinate sum, M-convexity means the following exchange condition. For all \(\alpha,\beta\in S\) and every \(i\) with \(\alpha_i>\beta_i\), there is a \(j\) with \(\alpha_j<\beta_j\) such that both \(\alpha-e_i+e_j\) and \(\beta+e_i-e_j\) belong to \(S\). Here \(e_i\) is the \(i\)-th unit vector.

\subsection{The central coefficient need not dominate repeated columns}\label{app:additional-2}

Consider

\[
A=\frac14\begin{pmatrix}1&1&2\\0&2&2\\3&1&0\end{pmatrix}.
\]

This matrix is doubly stochastic. Direct permanent sums give

\[
\per(A)=\frac5{16},\qquad
\per(A_{[:,(1,3,3)]})=\frac38>\frac5{16}.
\]

The matrix \(A_{[:,(1,3,3)]}\) repeats column three and omits column two. For a general multiplicity vector \(\alpha\), the repeated-column permanent equals \(\alpha!p_\alpha\). Therefore \(\alpha!p_\alpha\le p_{\mathbf1}\) is false. Summing this pointwise comparison cannot prove \eqref{eq:additional-1}. The actual strong deficit is positive:

\[
\per(AA^{\T})=\frac{17}{64},\qquad
\per(A)^2-c_3\per(AA^{\T})=\frac{89}{2304}.
\]

\subsection{Pairwise averaging need not decrease the strong ratio}\label{app:additional-3}

Write \(R(A)=\per(A)^2/\per(AA^{\T})\). For

\[
A=\frac12\begin{pmatrix}1&0&1\\1&1&0\\0&1&1\end{pmatrix},\qquad
B=\frac14\begin{pmatrix}1&1&2\\2&2&0\\1&1&2\end{pmatrix},
\]

the matrix \(B\) averages rows one and three of \(A\). Both matrices are doubly stochastic, and

\[
\per(A)=\per(B)=\frac14,
\quad\per(AA^{\T})=\frac14,
\quad\per(BB^{\T})=\frac{15}{64}.
\]

Hence \(R(A)=1/4<4/15=R(B)\).

A separate column example is

\[
C=\frac15\begin{pmatrix}2&3&0\\2&0&3\\1&2&2\end{pmatrix},\qquad
\widetilde C=\frac1{10}\begin{pmatrix}5&5&0\\2&2&6\\3&3&4\end{pmatrix}.
\]

The second matrix averages columns one and two. Exact evaluation gives

\[
R(C)=\frac{1089}{3841}<\frac{169}{590}=R(\widetilde C).
\]

Thus arbitrary row or column averaging cannot be assumed to decrease \(R\) toward \(R(J_n)=c_n\). These examples do not exclude a selected sequence of averaging steps.

\subsection{Normalized coefficients are not multivariate Jensen-concave}\label{app:additional-4}

Let

\[
A=\frac15\begin{pmatrix}2&0&3\\0&3&2\\3&2&0\end{pmatrix},\qquad
125p_A=(2x+3z)(3y+2z)(3x+2y).
\]

Put \(D_\alpha=\alpha!p_\alpha\). Expansion gives

\[
D_{111}=\frac{35}{125},\qquad
D_{210}=D_{021}=D_{102}=\frac{36}{125}.
\]

The average of the three latter exponent vectors is \((1,1,1)\). Nevertheless,

\[
\log D_{111}<\frac13(\log D_{210}+\log D_{021}+\log D_{102}).
\]

The polynomial is a balanced stable product of nonnegative linear forms. These properties do not give full multivariate Jensen concavity. The three weights in this comparison are \(1/3\). They are not the actual occupancy probabilities, so this example alone does not refute an entropy inequality with those probabilities.

\subsection{The collision polynomial can lose stability and balance}\label{app:additional-5}

For a product polynomial, define

\[
h_A(x)=\sum_\alpha\alpha!p_\alpha^2x^\alpha
=\per(A\operatorname{diag}(x)A^{\T}).
\]

Use the matrix \(C\) in Section~\ref{app:additional-3}. After multiplication by \(5^6\), its collision polynomial is

\[
\begin{aligned}
h(x,y,z)={}&96x^3+392x^2y+392x^2z+288xy^2\\
&+1089xyz+288xz^2+648y^2z+648yz^2.
\end{aligned}
\]

Its mass is \(3841\). Its three first-moment numerators are \(3521,4001,4001\). Thus its normalized mean is

\[
\left(\frac{3521}{3841},\frac{4001}{3841},\frac{4001}{3841}\right)
\ne(1,1,1).
\]

Also,

\[
\operatorname{Hess}(\partial_yh)=
\begin{pmatrix}784&576&1089\\576&0&1296\\1089&1296&1296\end{pmatrix}
\]

has determinant \(-120932352\) and trace \(2080\). It has two positive eigenvalues and one negative eigenvalue. The quadratic derivative therefore violates the Lorentzian Hessian condition. Differentiation preserves stability, so \(h\) is neither stable nor Lorentzian.

The same example also excludes an arbitrary unscaled singleton deletion. Terms with \(x\)-exponent one have total collision weight \(1665\). Therefore

\[
\frac{N([x]p_C)}{N(p_C)}=\frac{1665}{3841}<\frac49,
\qquad 14985<15364.
\]

This last calculation concerns the unscaled norm. It does not determine the norm after capacity balancing.

For a nonnegative matrix \(A\) with \(N(p_A)>0\), the normalized collision law has a useful exact mixture representation. Choose a permutation \(\sigma\) with probability proportional to \(\prod_i(AA^{\T})_{i,\sigma(i)}\). Conditional on this permutation, choose labels independently, with row-\(i\) probabilities proportional to \(a_{ij}a_{\sigma(i),j}\). Group the labels by occupancy. The resulting occupancy has mass \(\alpha!p_\alpha^2/N(p_A)\). Permutations with zero weight can be omitted. The conditional row distributions need not form a doubly stochastic matrix. Thus van der Waerden's theorem does not apply to each component merely from this representation.

\subsection{A global exponential trace lower bound fails}\label{app:additional-6}

The bound

\begin{equation}
\frac{\per(A)}{c_n}
\ge\exp\!\left(\frac12\|A-J_n\|_F^2\right)
\label{eq:additional-2}
\end{equation}

fails for

\[
A=\frac12\begin{pmatrix}0&1&1\\1&0&1\\1&1&0\end{pmatrix}.
\]

Here \(\per(A)/c_3=9/8\) and \(\|A-J_3\|_F^2=1/2\), whereas \(e^{1/4}>5/4>9/8\).

More generally, take \(A=(nJ_n-I)/(n-1)\). If \(D_n\) is the number of derangements of \(n\) elements, then

\[
\frac{\per(A)}{c_n}
=\frac{D_n n^n}{n!(n-1)^n},\qquad
\|A-J_n\|_F^2=\frac1{n-1}.
\]

At \(n=5\), the ratio is \(6875/6144\), but \(e^{1/8}>9/8>6875/6144\). These failures exclude \eqref{eq:additional-2}, not weaker trace bounds.

\subsection{Fourier magnitudes and a fixed transfer gap are insufficient}\label{app:additional-7}

For a circulant \(C(w)_{ij}=w_{j-i\bmod n}\), its Gram matrix depends only on its cyclic autocorrelation. At order three, let

\[
w=(4,1,1),\qquad v=(3,3,0).
\]

Both autocorrelations are \((18,9,9)\), but

\[
\per C(w)=78,\qquad
\per C(v)=54,\qquad
\per(C(w)C(w)^{\T})=11664.
\]

These values follow from \(\per C(a,b,c)=a^3+b^3+c^3+3abc\). After division by six, \(C(w)\) has eigenvalues \(1,1/2,1/2\). The nontrivial eigenvalues of \(C(v)/6\) have magnitudes \(1/2\) and phases \(\pm\pi/3\). Replacing phases by zero can therefore increase the permanent. A fixed-Gram argument must retain phase information.

A separate limitation follows at \(J_n\):

\[
\frac{\per(J_nJ_n^{\T})}{\per(J_n)^2}
=\frac{n^n}{n!},\qquad
\left(\frac{n^n}{n!}\right)^{1/n}\longrightarrow e.
\]

Thus no bound

\[
\per(CC^{\T})\le b_nK^n\per(C)^2,
\qquad b_n^{1/n}\longrightarrow1,
\]

can cover every circulant with one fixed \(K<e\). A transfer bound valid for arbitrary support must allow this gap to close.

\subsection{Two stability failures for circulant coefficient constructions}\label{app:additional-8}

First take \(C=I+P_4\), with a cyclic orientation for which

\[
p_C=(x_0+x_1)(x_1+x_2)(x_2+x_3)(x_3+x_0).
\]

For \(F_C=\sum\alpha!p_\alpha^2x^\alpha\), expansion gives

\[
\partial_{x_0}^2F_C
=4(x_1x_2+x_1x_3+2x_2^2+x_2x_3).
\]

At \(x_1=x_3=1,x_2=t\), this is \(4(2t^2+2t+1)\), with roots \((-1\pm i)/2\). Its Hessian block is

\[
\begin{pmatrix}0&4&4\\4&16&4\\4&4&0\end{pmatrix},
\]

with eigenvalues \(-4,10-2\sqrt{17},10+2\sqrt{17}\). Thus factorial-weighted coefficient squaring does not preserve stability or Lorentzianity, even for a circulant product. Cyclic symmetry still gives balance of \(F_C\). In fact, weighted arithmetic--geometric mean gives \(F_C(x)\ge F_C(\mathbf1)\prod_jx_j\) for \(x_j>0\). Balance does not repair the failed stability condition.

Second, let \(P\) be a four-cycle and put

\[
R_0=I+P+P^2+P^3,\quad R_1=P+P^2+P^3,\quad
R_2=P^2+P^3,\quad R_3=P^3.
\]

Define the mixed gap coefficient

\[
f(y)=[z_0z_1z_2z_3]
\prod_{i=0}^3\left(\sum_{j=0}^3\sum_{r=0}^3z_r(R_r)_{ij}y_j\right).
\]

The Hessian of \(\partial_{y_0}^2f\) is

\[
H=\begin{pmatrix}
24&36&48&48\\36&48&56&56\\48&56&56&60\\48&56&60&48
\end{pmatrix}.
\]

Its leading \(3\times3\) block has determinant \(-384\) and positive trace. More explicitly, at

\[
y_3=0,\qquad(y_0,y_1,y_2)=(-30+t,38+t,-13+t),
\]

the quadratic is \(204(1+t^2)\). This has nonreal roots. Differentiation and real specialization would preserve stability or give zero. Hence this extracted polynomial is not stable. The example does not give a negative coefficient in a circulant sorted-gap deficit.

\subsection{Transpose averaging does not restore stability or collision balance}\label{app:additional-9}

Let

\[
A=\frac16\begin{pmatrix}3&2&1\\1&0&5\\2&4&0\end{pmatrix},\qquad
A_0=6A,\qquad U=p_{A_0}+p_{A_0^{\T}}.
\]

The Hessian of \(\partial_yU\) is

\[
H=\begin{pmatrix}96&36&168\\36&0&120\\168&120&120\end{pmatrix},
\qquad\det H=-86400,\quad\operatorname{tr}H=216.
\]

It has two positive eigenvalues. An explicit stability obstruction uses \(x=(-12,10,5)\) and \(e=(1,1,1)\):

\[
x^{\T}Hx=24,\quad x^{\T}He=0,\quad e^{\T}He=864,
\qquad\tfrac12(x+te)^{\T}H(x+te)=12+432t^2.
\]

At \(t=i/6\), every coordinate of \(x+te\) has positive imaginary part, and this quadratic vanishes. Thus \(U\) is neither stable nor Lorentzian.

If \(U=\sum r_\alpha x^\alpha\), its collision weights satisfy

\[
\sum_\alpha\alpha!r_\alpha^2=54000,\qquad
\frac{\sum_\alpha\alpha\,\alpha!r_\alpha^2}{54000}
=\left(\frac{117}{125},\frac{129}{125},\frac{129}{125}\right).
\]

Transpose averaging therefore does not restore balanced collision moments. The matrix itself satisfies Gyires with positive deficit

\[
4\per(A)^2-4c_3Q(A)=\frac{169}{486}.
\]

\subsection{An operator inequality stronger than the scalar target fails}\label{app:additional-10}

Let \(T=A^{\otimes n}\). Let \(E\) project orthogonally onto tensor basis vectors whose indices are all distinct. In the permutation basis, \(S=ETE\) has entries

\[
S_{\sigma,\tau}=\prod_{i=1}^n a_{\sigma(i),\tau(i)}.
\]

Every row and column sum is \(\per(A)\). If \(v\) is the unit uniform vector on this subspace, then \(Sv=S^{\T}v=\per(A)v\). For symmetric \(A\), the scalar target would follow from

\begin{equation}
S^2\succeq c_nE(A^2)^{\otimes n}E.
\label{eq:additional-3}
\end{equation}

Use the symmetric order-three matrix from Section~\ref{app:additional-6}. Order the basis as

\[
(123),\ (132),\ (213),\ (231),\ (312),\ (321).
\]

Put \(w=e_{123}+e_{132}\) and \(G=E(A^2)^{\otimes3}E\). Exact matrix products give

\[
\langle w,S^2w\rangle=\frac1{16},\qquad
\langle w,Gw\rangle=\frac5{16},\qquad
\langle w,(S^2-\tfrac29G)w\rangle=-\frac1{144}.
\]

Thus \eqref{eq:additional-3} fails even on a nonnegative vector. On the required uniform vector \(v\), the scalar strong deficit is instead \(1/144\). The scalar problem requires information specific to that vector.

\subsection{Positive-semidefinite completion does not give coefficientwise positivity}\label{app:additional-11}

For a doubly stochastic \(A\), set

\[
S_A=\begin{pmatrix}0&A\\A^{\T}&0\end{pmatrix},\qquad
B_t=\frac{I_{2n}+tS_A}{1+t},\quad0\le t\le1.
\]

The bound \(\|A\|_2\le1\) makes \(B_t\) positive semidefinite and doubly stochastic. Also,

\[
\per(I_{2n}+tS_A)
=\sum_{k=0}^n t^{2k}
\sum_{|I|=|J|=k}\per(A[I,J])^2.
\]

The identity entries fix the unused rows and columns. The two cross-block matchings produce the same minor permanent twice. Since \(0\preceq B_t^2\preceq B_t\), permanent monotonicity in the positive-semidefinite order gives

\[
\per((I_{2n}+tS_A)^2)
\le(1+t)^{2n}\per(I_{2n}+tS_A).
\]

This valid inequality on \([0,1]\) does not imply coefficientwise positivity. The coefficient of \(t^{4n}\) in its right side minus its left side is

\[
\per(A)^2
-\per(AA^{\T})\per(A^{\T}A).
\]

For \(A=(I+P_4)/2\), it equals

\[
\frac1{64}-\left(\frac9{64}\right)^2=-\frac{17}{4096}.
\]

The same matrix gives a distinct limit for duplication:

\[
\per(A)=\frac18,\qquad
\per(AA^{\T})=\frac9{64},\qquad
\per(A\otimes\mathbf1_{2\times2})=\frac98,
\]

\[
\per(A)^2-
\frac{\per(A\otimes\mathbf1_{2\times2})}{\binom84}
=-\frac1{2240}.
\]

Thus the rank-two duplication bound cannot extend to all nonnegative doubly stochastic matrices. This example is also stated in Remark 1.3 of Kunisky, Spielman, and Yu, \emph{Inequalities for rank-two permanents and finite free convolutions}, \href{https://arxiv.org/html/2608.28520v1}{arXiv:2608.28520v1}. The strong deficit for this matrix is \(5/2048>0\).

\subsection{Matrix symmetrization and signed contractions do not give the required reductions}\label{app:additional-12}

For

\[
A=\frac15\begin{pmatrix}0&0&5\\1&4&0\\4&1&0\end{pmatrix},\qquad
B=\frac{A+A^{\T}}{2},
\]

exact evaluation gives

\[
\frac{\per(A)^2}{Q(A)}=\frac{289}{321},\qquad
\frac{\per(B)^2}{Q(B)}=\frac{110889}{122956}.
\]

The first ratio minus the second is \(-61085/39468876<0\). Plain symmetrization therefore need not decrease the original Gyires ratio.

The nonnegative-entry hypothesis also cannot be replaced by symmetry, contraction, and the van der Waerden lower bound. The signed matrix

\[
C=\frac16\begin{pmatrix}-2&4&4\\4&1&1\\4&1&1\end{pmatrix}
\]

is symmetric, has row and column sums one, and has eigenvalues \(1,-1,0\). Yet

\[
\per(C)=\frac5{18}>c_3,\qquad
\per(C^2)=\frac12,\qquad
\per(C)^2-c_3\per(C^2)=-\frac{11}{324}.
\]

In the centered rank-two representation \(C_{ij}=(1+x_iy_j)/3\), take \(x=(\sqrt2,-1/\sqrt2,-1/\sqrt2)\) and \(y=-x\). Then the normalized permanent is \(5/4\), and the normalized Gram permanent is \(9/4\). Centering and a singular-value bound do not replace entrywise nonnegativity.

\subsection{The occupancy entropy bridge fails with the actual probabilities}\label{app:additional-13}

For doubly stochastic \(A\), the coefficients \(p_\alpha\) of \(p_A\) form a probability law. Put

\[
q_\alpha=\frac{c_n}{\alpha!},\qquad r_\alpha=\frac{p_\alpha}{q_\alpha},\qquad
D(p\Vert q)=\sum_{p_\alpha>0}p_\alpha\log r_\alpha.
\]

The strong target is \(E_qr^2\le r_{\mathbf1}^2\). The proposed auxiliary inequality \(D(p\Vert q)\le\log r_{\mathbf1}\) fails for \(A=(I+P_n)/2\) and every \(n\ge4\).

To compute its exact gap, use independent fair bits \(X_1,\ldots,X_n\). The occupancy is \(\alpha_j=1-X_j+X_{j-1}\), with cyclic indices. Every nonconstant binary word is determined by its cyclic differences. The two constant words both give \(\mathbf1\). Hence

\[
p_{\mathbf1}=2^{1-n},\qquad p_\alpha=2^{-n}
\quad\text{for every other supported occupancy}.
\]

If \(K\) counts transitions of either fixed orientation around the cycle, then \(\alpha!=2^K\) and \(EK=n/4\). The entropy is

\[
-\sum_{p_\alpha>0}p_\alpha\log p_\alpha=(n-2^{1-n})\log2.
\]

Therefore

\[
D(p\Vert q)=-\frac{3n}{4}\log2+2^{1-n}\log2-\log c_n,
\]

\[
D(p\Vert q)-\log r_{\mathbf1}
=\left(\frac n4-1+2^{1-n}\right)\log2>0\quad(n\ge4).
\]

At \(n=4\), \(r_{\mathbf1}=4/3\) and \(D(p\Vert q)=\log(4/3)+(\log2)/8\). Still, \(E_qr^2=3/2<16/9=r_{\mathbf1}^2\). This failure concerns the entropy bridge. It makes no claim about a separate proposed comparison between order-two Rényi divergence and twice the relative entropy.

\subsection{Balance and negative dependence do not replace product structure}\label{app:additional-14}

Balance alone is insufficient. The law assigning mass \(1/n\) to each pure occupancy \(ne_j\), for \(n\ge2\), has mean \(\mathbf1\) and zero central mass. Its polynomial is \(\sum_jx_j^n/n\), with \(N=n!/n>0\). Thus its central coefficient cannot satisfy the norm inequality.

A stronger example also has negative association and bounded moments. For \(0\le\theta\le1\), let \(p^{(\theta)}\) put mass \(\theta\) at \((1,1,1)\) and mass \((1-\theta)/6\) at each permutation of \((2,1,0)\). Its mean is \(\mathbf1\), and

\[
\operatorname{Cov}(p^{(\theta)})=(1-\theta)(I_3-J_3)
\preceq I_3-J_3=\operatorname{Cov}(q).
\]

Every mixed falling-factorial moment is at most one. The first moments equal one. At degree two the moments are \(2(1-\theta)/3\) for a repeated coordinate and \((2+\theta)/3\) for distinct coordinates. At degree three they are \((1-\theta)/3\) for pattern \((2,1,0)\), \(\theta\) for \((1,1,1)\), and zero for \((3,0,0)\). Higher moments vanish.

To check negative association, condition on \(\alpha_1\). For values zero and two, the other pair is uniform on \(\{(1,2),(2,1)\}\) and \(\{(0,1),(1,0)\}\), respectively. At value one, it is a mixture of the uniform law on \(\{(0,2),(2,0)\}\) and the point mass at \((1,1)\). The upper conditional law dominates the middle law coordinatewise, and the middle dominates the lower law. This follows by coupling the two unequal points with the corresponding upper or lower points, and coupling \((1,1)\) with an equal mixture of those points. Thus every increasing function of the other coordinates has conditional expectation decreasing in \(\alpha_1\). Symmetry gives the same property for each coordinate. In three coordinates, disjoint nonempty coordinate sets always have a singleton side. This proves negative association.

At \(\theta=3/13\),

\[
p_{111}=\frac3{13}>c_3,\qquad
N(p)=\theta^2+\frac{(1-\theta)^2}{3}=\frac{127}{507},
\qquad c_3N(p)-p_{111}^2=\frac{11}{4563}>0.
\]

This is an abstract occupancy law, not a doubly stochastic product polynomial. For its generating polynomial, the Hessian of the \(x\)-derivative has eigenvalue \((1-4\theta)/3=1/39\) in the antisymmetric \((y,z)\) direction. Its complementary block has another positive eigenvalue. The polynomial is neither stable nor Lorentzian. The example excludes a replacement of product structure by just the stated dependence and moment conditions.

\subsection{A mixed-coefficient estimate fails for Lorentzian cubics}\label{app:additional-15}

For a balanced ternary cubic, write \(b=[xyz]p\), \(U_i=[x_i^3]p\), and

\[
[x_j]p=a_jx_k^2+b x_kx_\ell+c_jx_\ell^2.
\]

The squared Fischer norm after capacity balancing of this binary quadratic is \(H_j=b^2+4a_jc_j\). Direct expansion gives

\[
G(p):=3\sum_jH_j-4N(p)
=5b^2-\sum_j\bigl((a_j+c_j)^2+7(a_j-c_j)^2\bigr)-24\sum_iU_i^2.
\]

The separate proposed estimate

\begin{equation}
\sum_j\bigl((a_j+c_j)^2+7(a_j-c_j)^2\bigr)\le3b^2
\label{eq:additional-4}
\end{equation}

is false. Take

\[
p=(y+z)\left[\frac{(y+z)^2}{4\sqrt3}
+\frac{x(y+z)}2+\frac{x^2}{\sqrt3}\right].
\]

It is balanced, \(b=1\), \(U_x=0\), and \(p(\mathbf1)=2+4\sqrt3/3\). Its \(y\)- and \(z\)-derivatives are squares of real linear forms. Its \(x\)-derivative is \((y+z)((y+z)/2+2x/\sqrt3)\). The support is the cubic simplex with only \((3,0,0)\) removed, which is M-convex. These quadratic derivatives have at most one positive Hessian eigenvalue, so \(p\) is Lorentzian.

The exact values are

\[
N(p)=\frac{13}{3},\quad H_x=H_y=H_z=2,\quad
G(p)=\frac23,\quad
\sum_j\bigl((a_j+c_j)^2+7(a_j-c_j)^2\bigr)=\frac{10}{3}>3.
\]

The desired average-deletion inequality still holds. The quadratic factor in variables \(x,s=y+z\) has discriminant \(-1/12\). Therefore \(p\) is not real stable. This example does not refute the stable-only form of \eqref{eq:additional-4}.

\subsection{Three limits of scalar likelihood reductions}\label{app:additional-16}

Let \(P_x\) be the sum of independent Bernoulli variables with parameters \(x\in[0,1]^n\). Fix \(\sum_i x_i=t\), with \(0<t<n\), and let \(B=\operatorname{Bin}(n,t/n)\). A likelihood ratio is \(r(k)=P_x(k)/B(k)\).

\textbf{Noninteger interpolation.} Take \(n=3,t=7/4\) and \(P=\operatorname{Bin}(2,7/8)\), with one additional zero Bernoulli parameter. Relative to \(B=\operatorname{Bin}(3,7/12)\),

\[
r=(27/125,18/25,9/5,0),\qquad
E_Pr=\frac{1539}{1000}>
\frac14r(1)+\frac34r(2)=\frac{153}{100}.
\]

Thus linear interpolation at a noninteger mean does not preserve the integer-mean collision-norm bound.

\textbf{Relative log-concavity after noise.} Take

\[
n=3,\quad t=5/4,\quad
x=(5/8,5/8,0),\quad y=(1,1/8,1/8),\quad\theta=t/n.
\]

The binomial-noise parameter is \(\kappa_y=\sum_j(y_j-\theta)^2/[n\theta(1-\theta)]=7/10\). Set \(x'_i=\theta+\kappa_y(x_i-\theta)\). Then

\[
x'=(9/16,9/16,1/8),\qquad
P_{x'}=(343,931,693,81)/2048,\qquad
P_y=(0,49/64,7/32,1/64).
\]

At indices \(1,2,3\), the ratio \(P_y/P_{x'}\) is \((1568/931,64/99,32/81)\). It is not log-concave, because

\[
\frac{(64/99)^2}{(1568/931)(32/81)}=\frac{76}{121}<1.
\]

\textbf{Mean-interval endpoints.} Let \(T_\kappa\) be binomial noise on the count of \(n\) independent Bernoulli(\(\theta\)) bits. Each bit is retained with probability \(\kappa\); otherwise it is replaced by an independent Bernoulli(\(\theta\)) bit. Write \(\langle f,g\rangle_B=\sum B(k)f(k)g(k)\). For \(n=3\), take both laws to be \(P_t=1+\operatorname{Bern}(t-1)\), with \(1\le t\le2\). Set

\[
r_t=P_t/B,\qquad
\kappa_t=1-\frac{\operatorname{Var}(P_t)}{3\theta(1-\theta)},\qquad
D_t=\sqrt{\langle r_t,T_{\kappa_t}r_t\rangle_B}
-\langle r_t,r_t\rangle_B.
\]

At each endpoint, \(D_t=\sqrt{9/4}-9/4=-3/4\). At \(t=3/2\),

\[
r_t=(0,4/3,4/3,0),\qquad D_t=\sqrt{31/27}-4/3>-3/4.
\]

Thus the maximum cannot in general be restricted to the two endpoints. This example does not violate the desired inequality \(D_t\le0\).

\subsection{A collision probability need not decrease with the trial counts}\label{app:additional-17}

A more general collision comparison has two proved boundary cases. Let \(r,s\) be positive integers, set \(v=rs/(r+s)\), and let \(h\ge r\), \(N\ge s\) be integers. The proposed bound is

\begin{equation}
\Pr\{\operatorname{Bin}(h,v/h)=\operatorname{Bin}(N,v/N)\}
\le\binom{r+s}{r}\left(\frac r{r+s}\right)^r
\left(\frac s{r+s}\right)^s.
\label{eq:additional-6}
\end{equation}

The two binomial variables are independent. The right side is the collision probability at \(h=r,N=s\). Bound \eqref{eq:additional-6} holds on both edges \(h=r\) and \(N=s\).

To prove this, put \(\gamma=s/(r+s)\). Take independent \(R\sim\operatorname{Bin}(h,r/h)\) and \(S\sim\operatorname{Bin}(N,s/N)\). Retain each success of \(R\) with probability \(\gamma\), and each success of \(S\) with probability \(1-\gamma\). The retained counts have the laws on the left of \eqref{eq:additional-6}. Given \(R,S\), their collision probability is

\[
H(R,S)=\binom{R+S}{R}(1-\gamma)^R\gamma^S.
\]

This identity follows by summing the two binomial probabilities and using Vandermonde's identity. If \(h=r\), then \(R=r\) deterministically, and

\[
\frac{H(r,S+1)}{H(r,S)}
=\gamma\frac{r+S+1}{S+1}.
\]

The maximum is attained at \(S=s-1,s\), so averaging proves \eqref{eq:additional-6}. Interchanging the two sides proves the other edge. The interior case is not proved by this argument.

For \(r=s=1\), the endpoint collision comparison asks whether

\[
\Pr\{\operatorname{Bin}(h,1/(2h))
=\operatorname{Bin}(N,1/(2N))\}\le\frac12,
\qquad h,N\ge1,
\]

where the two variables are independent. A stronger assertion that the probability decreases in each trial count is false. With equal counts, the probability is

\[
L_h=\sum_{k=0}^h\binom hk^2
\left(\frac1{2h}\right)^{2k}
\left(1-\frac1{2h}\right)^{2h-2k}.
\]

Exact evaluation gives

\[
L_2=\frac{59}{128}<\frac{1154188663}{2500000000}=L_5.
\]

The limiting value is \(e^{-1}\sum_{k\ge0}1/[4^k(k!)^2]\), also written \(e^{-1}I_0(1)\). Both the finite increase and this limit concern the monotonicity assertion. They do not refute the bound \(1/2\).

\subsection{Column averaging can increase a rank-two permanent}\label{app:additional-18}

Let

\[
A=\frac1{30}\begin{pmatrix}4&15&11\\4&15&11\\22&0&8\end{pmatrix},\qquad
\widetilde A=\frac1{30}\begin{pmatrix}4&13&13\\4&13&13\\22&4&4\end{pmatrix}.
\]

Both matrices are doubly stochastic and have rank two. The second averages columns two and three, but

\[
\per(A)=\frac{137}{450},\qquad
\per(\widetilde A)=\frac{689}{2250},\qquad
\per(\widetilde A)-\per(A)=\frac2{1125}>0.
\]

Hence arbitrary column mixing does not decrease the permanent, even in rank two. The same matrix also satisfies

\[
\per(A_{[:,(1,2,2)]})=\frac{11}{30}>\frac{137}{450}.
\]

Thus the scalar likelihood-peak theorem does not give a pointwise maximum at the central multivariate occupancy, even in rank two.

\subsection{Minimum Fischer norm cannot supply a proposed deletion factor}\label{app:additional-19}

For a nonnegative degree-\(n\) polynomial in \(n\) variables, define

\[
\mathcal C_F(p)=\inf_{D>0}\frac{N(p(Dx))}{(\det D)^2},\qquad
g_n=\left(\frac{n-1}{n}\right)^{n-1},\quad g_1=1.
\]

The proposed comparison

\begin{equation}
\frac{\mathcal C_F(p_A)}{N(p_A)}\ge\frac{g_n}{g_d},
\label{eq:additional-5}
\end{equation}

where \(d\) is the number of positive entries in a zero-containing column, fails even in rank two.

Take \(n=50\),

\[
u=(0,1/2,\underbrace{1/96,\ldots,1/96}_{48}),\qquad
A_{i1}=u_i,\quad A_{ij}=(1-u_i)/49\quad(j\ge2).
\]

This doubly stochastic matrix has rank two and exactly one zero entry. Thus \(d=49\). Let

\[
P=\operatorname{Bern}(1/2)+\operatorname{Bin}(48,1/96),\qquad
B=\operatorname{Bin}(50,1/50),\qquad
C(z)=\sum_{k=0}^{49}\frac{P_k^2}{B_k}z^k.
\]

Writing \(Y=x_2+\cdots+x_{50}\) gives \(p_A=\sum_kP_kx_1^k(Y/49)^{50-k}\). The multinomial theorem then gives \(N(p_A)=c_{50}C(1)\). Scaling only \(x_1\) by \(\sqrt z\) has squared determinant \(z\), so

\[
\frac{\mathcal C_F(p_A)}{N(p_A)}\le\frac{C(z)}{zC(1)}.
\]

For a direct rational certificate, put

\[
P_k=\frac{\binom{48}{k}95^{48-k}
+\binom{48}{k-1}95^{49-k}}{2\cdot96^{48}},\qquad
B_k=\frac{\binom{50}{k}49^{50-k}}{50^{50}},
\]

with each out-of-range binomial term set to zero before evaluation. Substitution and rational comparison give

\[
\frac{C(31/32)}{(31/32)C(1)}
<\frac{99973}{100000}<\frac{99979}{100000}
<\frac{g_{50}}{g_{49}}.
\]

This proves the failure of \eqref{eq:additional-5} without an optimizing scale.

The failure persists in all sufficiently large dimensions of the same family. Fix \(0<a<1\), write \(b=1-a\), and replace \(u\) by \((0,a,b/(n-2),\ldots,b/(n-2))\). The occupancy law tends to \(\operatorname{Bern}(a)+\operatorname{Pois}(b)\), while the reference tends to \(\operatorname{Pois}(1)\). The limiting collision series is

\[
C_\infty(z)=e^{1-2b+b^2z}\bigl[b^2(1+az)^2+a^2z\bigr].
\]

For \(1\le k\le n\), the finite laws satisfy

\[
P_{n,k}\le\frac{b^{k-1}(b^2+ak)}{k!},\qquad
B_{n,k}^{-1}\le e^{k+2}k!.
\]

Bound the \(k=0\) term separately. Set the finite collision coefficients to zero for \(k>n\). These estimates give summable factorial majorants for the collision series and its derivative on compact \(z\)-intervals. Dominated convergence applies. Direct differentiation gives

\[
\frac{C_\infty'(1)}{C_\infty(1)}-1
=\frac{a^4b^2}{1-a^2+a^4}>0.
\]

Thus a fixed scale just below one decreases \(C_\infty(z)/z\), while \(g_n/g_{n-1}\to1\). This is the stated dimension-persistent failure. It does not refute the capacity-balanced deletion inequality, the strong Gram inequality, or Gyires's inequality.

\subsection{An entropy reduction for scalar reference laws}\label{app:additional-20}

Let \(n,m\) be integers with \(1\le m\le n-1\), and put \(B=\operatorname{Bin}(n,m/n)\). Let \(P\) have mean \(m\), and suppose \(r=P/B\) is log-concave with interval support. Define

\[
D(P\Vert B)=\sum_kP_k\log r_k,\qquad
D_2(P\Vert B)=\log\sum_kB_kr_k^2.
\]

Concavity of the piecewise-linear interpolation of \(\log r\) and the mean condition give

\[
D(P\Vert B)\le\log r_m.
\]

Consequently \(D_2(P\Vert B)\le2D(P\Vert B)\) would imply \(\sum_kB_kr_k^2\le r_m^2\). On a localized interval law, take integers \(0\le a\le m\le b\le n\) and \(t>0\). Set \(r_k=t^{k-m}/Z\) for \(a\le k\le b\), and set it to zero otherwise. Here \(Z=\sum_{k=a}^bB_kt^{k-m}\). With the tilt chosen to give mean \(m\), one has \(D(P\Vert B)=-\log Z=\log r_m\). Thus the entropy comparison is equivalent to the norm comparison on those localized laws. This reduction does not prove the comparison for asymmetric integer means outside the established cases.

An induction by division of a common variable has a separate failure in the Lorentzian class. The binary cubic

\[
p(X,Y)=Y^3+3XY^2+3X^2Y
\]

is Lorentzian because its coefficients divided by \(\binom3k\) are \((1,1,1,0)\), a log-concave sequence with interval support. But

\[
p/Y=Y^2+3XY+3X^2
\]

has discriminant \(-3\), so it is not a Lorentzian binary quadratic. Division by a common variable therefore does not preserve Lorentzianity in general. This failure does not apply to division within the real-stable class.

Two proper truncations give concrete norm ratios. For \(B=\operatorname{Bin}(n,1/n)\), set \(u_k=t^{k-1}\) for \(0\le k\le L\), and zero otherwise. Choose \(t>0\) so that \(\sum_k(k-1)B_ku_k=0\). Normalizing \(B_ku_k\) gives a mean-one law. Its scalar norm ratio is

\[
\frac{r_1^2}{\sum_kB_kr_k^2}=\frac1{\sum_kB_ku_k^2}.
\]

For \(n=3,L=2\), the mean equation gives \(t^2=4/3\), and

\[
\sum_kB_ku_k^2=\frac29+\frac49+\frac8{27}=\frac{26}{27}.
\]

The ratio is exactly \(27/26\). For \(n=4,L=3\), the tilt is the unique positive root of \(8t^3+18t^2-27=0\). The ratio is

\[
\frac{256}{81/t^2+108+54t^2+12t^4}
\approx1.003921142546.
\]

This value is below \(256/255\approx1.003921568627451\), but remains above one. The displayed root equation specifies this example exactly. The all-order scalar proof does not depend on this numerical comparison.

\subsection{An equivalent product inequality}\label{app:additional-21}

In each fixed order, universal validity of the strong inequality \eqref{eq:additional-1} is equivalent to

\begin{equation}
c_n\per(AB)
\le\per(A)\per(B)
\label{eq:additional-7}
\end{equation}

for all doubly stochastic \(A,B\). Indeed, permanent Cauchy--Schwarz gives

\[
\per(AB)^2
\le\per(AA^{\T})\per(B^{\T}B).
\]

Apply \eqref{eq:additional-1} to \(A\) and \(B^{\T}\), then take nonnegative square roots to obtain \eqref{eq:additional-7}. Conversely, use \(B=A^{\T}\) in \eqref{eq:additional-7}. Thus known special families do not automatically give closure under matrix multiplication. That closure would need an additional permanent comparison.

\subsection{Tensor forms of the two inequalities}\label{app:additional-22}

Let \(e_1,\ldots,e_n\) be the standard orthonormal basis, and define

\[
u=n^{-1/2}(1,\ldots,1),\qquad
v=\frac1{\sqrt{n!}}\sum_{\sigma\in S_n}
e_{\sigma(1)}\otimes\cdots\otimes e_{\sigma(n)}.
\]

Then \(\|v\|=1\) and \(\langle v,u^{\otimes n}\rangle=\sqrt{c_n}\). Expansion of the permutation sums gives

\[
\per(A)=\langle v,A^{\otimes n}v\rangle,\qquad
\per(AA^{\T})=\|(A^{\T})^{\otimes n}v\|^2.
\]

For \(h=(A^{\T})^{\otimes n}v\), double stochasticity implies \(\langle h,u^{\otimes n}\rangle=\sqrt{c_n}\). Thus \eqref{eq:additional-1} says that the cosine of the angle between \(h\) and \(v\) is at least \(\sqrt{c_n}\). Section~\ref{app:additional-14} shows why balanced nonnegative coefficient laws alone cannot establish the needed norm bound.

For the original inequality, put \(T=A^{\otimes n}\) and \(H=(T+T^*)/2\), where \(T^*\) is the transpose operator. Direct multiplication gives

\[
\langle v,Hv\rangle=\per(A),\qquad
\|Hv\|^2=Q(A).
\]

The original inequality is therefore \(\langle v,Hv\rangle^2\ge c_n\|Hv\|^2\). This identity retains both matrix orientations. It is a reformulation, not a proof of the angle bound.

\subsection{The first two gap layers for Latin-square weight matrices}\label{app:additional-23}

Let \(n\ge3\), and let every row and column of \(B\) be a permutation of a vector \(w\). Let \(E\) be the all-ones matrix, put \(s=\sum_iw_i\), and define

\[
F(A)=n^n\per(A)^2-n!\per(AA^{\T}),\qquad
V(w)=\sum_{i<j}(w_i-w_j)^2,
\]

\[
W(w)=\sum_{i<j}(w_i-w_j)^2(w_i+w_j).
\]

The first two possibly nonzero powers in the expansion about \(tE\) are

\begin{equation}
\frac{F(tE+B)}{n^n n!(n-2)!}
=t^{2n-2}V(w)+t^{2n-3}H_3(w)
+\text{terms of degree at most }2n-4\text{ in }t,
\label{eq:additional-8}
\end{equation}

where

\begin{equation}
H_3(w)=\left(\frac{2(n-1)}n-\frac8{3n(n-2)}\right)sV(w)
+\frac4{3(n-2)}W(w).
\label{eq:additional-9}
\end{equation}

Here is a derivation. Set \(X=B-(s/n)E\) and \(a=t+s/n\). Then \(X\) has zero row and column sums. Expanding the permanent by the entries selected from \(X\) gives

\[
\per(aE+X)=n!a^n
+\frac{(n-2)!}{2}a^{n-2}\sum_{ij}X_{ij}^2
+\frac{2(n-3)!}{3}a^{n-3}\sum_{ij}X_{ij}^3
+\cdots.
\]

The quadratic and cubic coefficients follow by summing products over distinct rows and columns and then using the zero sums. For the cubic coefficient, inclusion--exclusion for three row labels and three column labels leaves \(4\sum_{ij}X_{ij}^3\) in the ordered sum; division by \(3!\) gives \(2\sum X_{ij}^3/3\).

Also, \((aE+X)(aE+X)^{\T}=na^2E+XX^{\T}\). The term linear in \(XX^{\T}\) in its permanent vanishes, since the entries of \(XX^{\T}\) sum to zero. Thus the Gram permanent has no contribution of degree two or three in \(X\). The constant terms cancel in \(F\). After the normalization in \eqref{eq:additional-8}, the next terms are

\[
a^{2n-2}\sum X_{ij}^2
+\frac4{3(n-2)}a^{2n-3}\sum X_{ij}^3.
\]

Every row of \(X\) permutes \(w-s/n\), so

\[
\sum X_{ij}^2=V(w),\qquad
\sum X_{ij}^3=W(w)-\frac{2s}{n}V(w).
\]

Substitute \(a=t+s/n\) and collect the two highest remaining powers of \(t\). This proves \eqref{eq:additional-8}--\eqref{eq:additional-9}.

For nonnegative sorted \(w\), write each coordinate as a sum of consecutive nonnegative gaps. Each squared difference, each sum \(w_i+w_j\), and \(s\) then has nonnegative coefficients. Both scalar factors in \eqref{eq:additional-9} are nonnegative for \(n\ge3\). Therefore the full quadratic and cubic gap layers have nonnegative coefficients, for every Latin-square arrangement. This statement gives no sign control over higher gap degrees.

\subsection{A two-product relaxation for balanced stable ternary cubics}\label{app:additional-24}

This reduction uses the definite determinantal representation theorem for ternary hyperbolic polynomials. A reference is Hanselka, \emph{Definite Determinantal Representations of Ternary Hyperbolic Forms}, Theorem 1, \href{https://arxiv.org/html/1411.1661v2}{arXiv:1411.1661v2}. Let \(p\) be a balanced stable ternary cubic with nonnegative coefficients and \(p(\mathbf1)=1\). The theorem gives a representation

\[
p(x,y,z)=\det(xA+yB+zC),\qquad A,B,C\succeq0,
\]

\[
A+B+C=I_3,\qquad
\operatorname{tr}A=\operatorname{tr}B=\operatorname{tr}C=1.
\]

Congruence normalizes the matrix at \(\mathbf1\) to the identity. Differentiating the determinant at the identity gives the traces. The positive orthant lies in the closure of the hyperbolicity cone, so the three coefficient matrices are positive semidefinite.

First assume \(I-A\) is positive definite. Congruence by \((I-A)^{-1/2}\), followed by an orthogonal change of basis, gives

\[
p=\det(I-A)\det\!\left(xR+
\operatorname{diag}(\beta_i y+(1-\beta_i)z)\right),
\quad0\le\beta_i\le1.
\]

Let \(r_1,r_2,r_3\) be the eigenvalues of \(R\), and let \(W\) contain the squared entries of an orthogonal eigenvector matrix. It is doubly stochastic. In dimension three,

\[
R_{ii}=\sum_jW_{ij}r_j,\qquad
(\operatorname{adj}R)_{ii}=\sum_jW_{ij}\prod_{k\ne j}r_k.
\]

Expansion against the diagonal matrix shows that the full coefficient vector is affine in \(W\). A Birkhoff decomposition therefore expresses \(p\) as a convex combination of

\begin{equation}
f_\pi(x,y,z)=\prod_{i=1}^3
\left(a_{\pi(i)}x+(1-a_{\pi(i)})\beta_i y
+(1-a_{\pi(i)})(1-\beta_i)z\right),
\label{eq:additional-10}
\end{equation}

where \(a_j=r_j/(1+r_j)\) are the eigenvalues of \(A\). Thus \(0\le a_j\le1\) and \(\sum_j a_j=1\). The factor \(\det(I-A)\) is absorbed by the three row normalizations.

Every \(f_\pi\) has value one and \(x\)-derivative one at \(\mathbf1\). Its \(y\)-derivative there is

\[
B_\pi=\sum_i(1-a_{\pi(i)})\beta_i.
\]

The mixture is balanced exactly when its average \(B_\pi\) is one. A hyperplane section of the order-three Birkhoff polytope has vertices at original vertices or at intersections with edges. Each edge joins two permutation matrices whose relative permutation is one cycle. Hence, for fixed \(a,\beta\), a sufficient convex-cone argument reduces to balanced mixtures

\[
(1-\theta)f_\pi+\theta f_\sigma,\qquad
(1-\theta)B_\pi+\theta B_\sigma=1,
\]

where \(\pi^{-1}\sigma\) is a transposition or a three-cycle. Original vertices are included by an endpoint value of \(\theta\).

This applies to the norm cone \(\{p:[xyz]p\ge\sqrt{c_3}\sqrt{N(p)}\}\), which is convex by the triangle inequality. It also applies to the average-deletion cone \(G(p)\ge0,[xyz]p\ge0\) from Section~\ref{app:additional-15}, whose defining formula is a Lorentz cone. The individual products in \eqref{eq:additional-10} need not be balanced. The relaxed two-product mixtures need not be stable. No assertion that all these mixtures belong to either target cone is proved here.

If \([x^3]p=0\), then \(\det A=0\), and the eigenvalues can be written \((t,1-t,0)\). If an eigenvalue of \(A\) equals one, positivity and \(A+B+C=I\) make the other coefficient matrices vanish on its eigenspace. They preserve the orthogonal complement, so the polynomial factors. Limits also cover this endpoint. The determinant reduction does not by itself settle the remaining boundary inequality.

\subsection{Two explicit constructions within polynomial deletion}\label{app:additional-25}

Let \(m_1,\ldots,m_n\) be nonnegative integers, and put \(M=\sum_i m_i\ge n\). Consider \(M\) factors, with \(m_i\) copies of \(x_i+y_1+\cdots+y_{M-n}\). Extract the coefficient of \(y_1\cdots y_{M-n}\). Choosing the \(n\) factors that supply an \(x\)-variable and assigning the distinct \(y\)-variables to the remaining factors gives

\[
[y_1\cdots y_{M-n}]
\prod_{i=1}^n(x_i+y_1+\cdots+y_{M-n})^{m_i}
=(M-n)!\,e_n(x_1^{[m_1]},\ldots,x_n^{[m_n]}).
\]

The notation \(x_i^{[m_i]}\) means that the argument list contains \(m_i\) copies of \(x_i\). This identity places the repeated-variable elementary-symmetric family inside the product-deletion class. It does not prove a norm inequality for that family in arbitrary order.

A second construction has a simpler deletion identity. Let \(n\ge2\) and

\[
p=(s+tx_1)^n-(tx_1)^n,\qquad
s=x_2+\cdots+x_n,\qquad t>0.
\]

Every singleton deletion is a positive scalar multiple of a positive linear form raised to power \(n-1\). Indeed,

\[
[x_1]p=nt\,s^{n-1},\qquad
[x_j]p=n(s-x_j+tx_1)^{n-1}\quad(j\ge2).
\]

Let \(b=[x_1\cdots x_n]p=n!t\). The central coefficient of each deletion is \(b\). A balanced pure power of degree \(n-1\) with central coefficient \(b\) has squared Fischer norm \(b^2/c_{n-1}\). Hence

\[
H([x_j]p)=\frac{b^2}{c_{n-1}},\qquad
\frac{H([x_j]p)}{g_nN(p)}=\frac{b^2}{c_nN(p)}
\quad\text{for every }j.
\]

These identities remain true after positive variable scaling, with the corresponding central coefficient and norm. Thus this family cannot separate the capacity-balanced deletion comparison from the final coefficient-norm comparison.

\subsection{A Hessian criterion and a family of Lorentzian coefficients}\label{app:additional-26}

Let \(H\) be real symmetric, let \(e=\mathbf1\), and suppose \(e^{\T}He>0\). The matrix \(H\) has at most one positive eigenvalue if and only if

\[
\frac{(He)(He)^{\T}}{e^{\T}He}-H\succeq0
\quad\text{on }e^\perp.
\]

To prove the equivalence, decompose a vector into a multiple of \(e\) and a vector in \(e^\perp\). Complete the square in the positive \(e\)-coordinate. The remaining quadratic form is the negative of the displayed restricted matrix. Sylvester's inertia law gives the assertion.

This criterion can be applied to a concrete coefficient family. For degree \(n\ge2\) in \(n\) variables, take

\[
p_\alpha=\frac{\lambda^{\sum_i\binom{\alpha_i}{2}}}{\alpha!},
\qquad0<\lambda<1,\quad|\alpha|=n.
\]

For \(|\beta|=n-2\), the Hessian of \(\partial^\beta p\) is a positive diagonal congruence, with a positive scalar factor, of \(E-(1-\lambda)I\), where \(E\) is the all-ones matrix. Its diagonal entries before congruence are \(\lambda\), and its off-diagonal entries are one. It has one positive eigenvalue and \(n-1\) negative eigenvalues. The support is the full degree-\(n\) simplex, so \(p\) is Lorentzian.

One can also set a selected pure coefficient \(p_{ne_j}\) to zero. This subtracts a nonnegative diagonal rank-one matrix from the Hessian for \(\beta=(n-2)e_j\), and leaves the other Hessians unchanged. It cannot increase their numbers of positive eigenvalues. The support is the degree simplex with that vertex removed, which is M-convex. The resulting polynomial is therefore Lorentzian as well. This construction supplies admissible examples. It gives no global optimum or norm bound.

\subsection{The two Gram permanents need not have a fixed order}\label{app:additional-27}

A matrix with three exceptional rows need not satisfy \(\per(A^{\T}A)\le\per(AA^{\T})\). Take

\[
A=\frac15\begin{pmatrix}
0&3&2&0&0\\
2&0&0&3&0\\
1&0&1&0&3\\
1&1&1&1&1\\
1&1&1&1&1
\end{pmatrix}.
\]

Every row and column sums to one. Its last two rows equal those of \(J_5\). Direct permanent sums give

\[
\per(A)=\frac{42}{625},\qquad
\per(AA^{\T})=\frac{19926}{390625},\qquad
\per(A^{\T}A)=\frac{508602}{9765625}.
\]

Thus

\[
\per(A^{\T}A)-\per(AA^{\T})
=\frac{10452}{9765625}>0.
\]

Both strong deficits are positive. They are \(624276/244140625\) for the row Gram matrix and \(15356052/6103515625\) for the column Gram matrix. Transposition reverses the displayed strict ordering. A proof for one Gram orientation therefore needs an additional argument for the other orientation.

\subsection{Two alternative bounds for the three-shift transfer matrices}\label{app:additional-28}

Let \(C=aP_n^{-1}+bI+cP_n\), where \(a,b,c\ge0\), and define

\[
L=\max\left\{a,c,\frac{b+\sqrt{b^2+4ac}}2\right\},\quad
 d=a^2+b^2+c^2,\quad e=b(a+c),\quad f=ac.
\]

The transfer identity for this family has two scalar blocks \(f\), two blocks with the spectrum of \(T_1\), and one block \(T_2\), where

\[
T_1=\begin{pmatrix}e&d&e&f\\f&0&0&0\\0&f&0&0\\0&0&f&0\end{pmatrix},
\qquad
T_2=\begin{pmatrix}
d&e&0&f&0&0\\e&0&e&0&f&0\\f&0&0&0&0&0\\
0&e&d&0&0&f\\0&f&0&0&0&0\\0&0&f&0&0&0
\end{pmatrix}.
\]

Consequently,

\[
\per(CC^{\T})=2f^n+2\operatorname{tr}(T_1^n)+\operatorname{tr}(T_2^n),
\qquad \per(C)\ge L^n.
\]

The case \(L=0\) is immediate. Otherwise divide the weights by \(L\). Then \(0\le a,c\le1\) and \(b+ac\le1\). Increasing \(b\) to \(1-ac\) increases both transfer matrices entrywise. The following arguments therefore give bounds for all the normalized weights.

\textbf{An explicit rational comparison vector.} Put \(K=51/20\), \(k=K-2=11/20\), and \(f=ac\). Since \(a+c\le1+f\),

\[
e\le1-f^2,\qquad d\le2-2f+2f^2,\qquad 0\le f\le1.
\]

For \(v=(1,k,K^{-1},k,kK^{-1},K^{-2})^{\T}\), the first row of \(T_2v\le Kv\) follows from

\[
d+k(e+f)\le K-(2-k)f(1-f)\le K.
\]

The second row follows from

\[
(1+K^{-1})e+(k/K)f
\le\frac{K+1}{K}+\frac{k^2}{4K(K+1)}<Kk.
\]

The final strict difference is \(2881/1448400\). The fourth row is at most

\[
k(1-f^2)+(2-2f+2f^2)/K+f/K^2.
\]

Its quadratic coefficient is \(239/1020>0\). Its maximum on \([0,1]\) is therefore at an endpoint. The endpoint differences from \(Kk\) are \(1391/20400\) and \(483161/1040400\), both positive. Rows three, five, and six use only \(f\le1\).

For \(u=(1,K^{-1},K^{-2},K^{-3})^{\T}\), the first row of \(T_1u\le Ku\) follows from

\[
1+2/K+1/K^2+1/K^3<K,
\]

with difference \(1463381/2653020\). The other rows again use \(f\le1\). For any nonnegative matrix \(T\) and positive vector \(v\) with \(Tv\le Kv\),

\[
(T^n)_{ii}v_i\le(T^nv)_i\le K^nv_i.
\]

After restoring the scale, these bounds give \(\per(CC^{\T})\le16K^nL^{2n}\). The exact inequalities

\[
94^{94}20^{94}>16\cdot94!\cdot51^{94},\qquad
20\cdot95^{94}>51\cdot94^{94}
\]

show that \(16c_nK^n<1\) for every \(n\ge94\). Indeed, \((1+1/n)^n\) increases with \(n\). This bound proves the tail \(n\ge94\). Its finite range is \(3\le n\le93\), together with the elementary orders one and two.

\textbf{A leading-minor certificate with constant \(5/2\).} With \(b=1-ac\), form \(5I-2T_1\) and \(5I-2T_2\). Every leading principal minor has positive coefficients in its tensor-product Bernstein basis on \([0,1]^2\). In increasing minor size, the smallest coefficients are

\[
T_1:\quad 3,\ 37/3,\ 65,\ 325,
\qquad
T_2:\quad 1,\ 1,\ 5,\ 25,\ 125,\ 609.
\]

These finite comparisons use exact rational arithmetic and require no subdivision. They can be reproduced directly from the displayed matrices. If a minor is \(p(a,c)=\sum p_{rs}a^rc^s\) with bidegree \((m,h)\), its Bernstein coefficient at \((i,j)\) is

\[
\beta_{ij}=\sum_{r\le i,\ s\le j}p_{rs}
\frac{\binom ir}{\binom mr}\frac{\binom js}{\binom hs}.
\]

The Bernstein basis functions are nonnegative and sum to one. Each minor is therefore strictly positive on the whole square.

The matrices \(5I-2T_j\) have nonpositive off-diagonal entries. Elimination without row swaps has positive pivots, since they are ratios of consecutive leading minors. Every Schur complement retains nonpositive off-diagonal entries. The triangular factors thus have nonnegative inverses. It follows that \((5I-2T_j)^{-1}\ge0\). Its positive diagonal gives

\[
x=(5I-2T_j)^{-1}\mathbf1>0,\qquad
T_jx=(5x-\mathbf1)/2<(5/2)x.
\]

The spectral radius is therefore less than \(5/2\). The full transfer bound is \(16(5/2)^nL^{2n}\). The exact comparisons

\[
70^{70}2^{70}>16\cdot70!\cdot5^{70},\qquad
\frac52\left(\frac67\right)^6=\frac{116640}{117649}<1
\]

give the infinite tail from order \(70\). Its finite range is \(3\le n\le69\), together with the elementary orders one and two. The sharper constant \(1+\sqrt2\) in the main theorem reduces the tail threshold to \(48\).

\textbf{The finite ranges.} An exact integer coefficient certificate covers all three reflection-reduced sorting chambers for every \(3\le n\le93\). It checks 273 chambers and 1,674,309 coefficient positions. No coefficient is negative. Reflection exchanges \(a\) and \(c\), so these chambers cover all six weight orders. The certificate uses characteristic-polynomial recurrences and Newton identities in sorted-gap coordinates. Run \texttt{node\ independent\_finite.mjs\ 93} to reproduce it. The program writes \nolinkurl{independent_finite.json} in the current directory. Compare that output with \nolinkurl{results/permutation-alternative-finite.json} in the supporting material. The reproduction guide gives commands that keep the two finite ranges in separate files.

The certificate supplies both finite ranges stated above. Combining it with either tail bound, the elementary orders one and two, and cycle decomposition gives a separate proof for all permutation polynomials \(aI+bQ+cQ^2\) with nonnegative weights. This conclusion concerns this three-term family. It does not establish the inequality for general circulants.
